\documentclass[12pt,a4paper]{article}
% ==============================================================================
% PREÁMBULO ESTANDARIZADO UNSCH - CIENCIAS DE LA SALUD (ENFERMERÍA)
% Compatible tanto con pdflatex como con xelatex (Unicode)
% ==============================================================================

\usepackage{iftex}

\ifPDFTeX
  \usepackage[utf8]{inputenc}
  \usepackage[T1]{fontenc}
  \usepackage{mathptmx} % Fuente Times 12pt requerida por la normativa
\else
  \usepackage{fontspec}
  \setmainfont{Times New Roman}
\fi

\usepackage[spanish,es-nodecimaldot,es-tabla]{babel}

% Márgenes normativos UNSCH:
% Superior: 3.0 cm | Izquierdo: 3.5 cm | Derecho: 2.5 cm | Inferior: 2.5 cm
\usepackage[top=3.0cm, left=3.5cm, right=2.5cm, bottom=2.5cm, a4paper]{geometry}

% Interlineado reglamentario (1.5 líneas)
\usepackage{setspace}
\onehalfspacing

% Sangría reglamentaria de párrafo
\setlength{\parindent}{1.27cm} % Sangría estilo APA (media pulgada)
\setlength{\parskip}{6pt}      % Separación ligera entre párrafos

% Paquetes matemáticos y símbolos
\usepackage{amsmath,amssymb,amsfonts}

% Tablas y matrices avanzadas
\usepackage{booktabs}
\usepackage{tabularx}
\usepackage{array}
\usepackage{multirow}
\usepackage{longtable}

% Inclusión de imágenes y gráficos
\usepackage{graphicx}
\usepackage{float}
\usepackage{caption}
\captionsetup{font=small, labelfont=bf, justification=centering}

% Encabezados y pie de página
\usepackage{fancyhdr}
\pagestyle{fancy}
\fancyhf{}
\fancyhead[R]{\small\thepage}
\fancyhead[L]{\small\nouppercase{\leftmark}}
\renewcommand{\headrulewidth}{0.5pt}
\renewcommand{\footrulewidth}{0pt}
\setlength{\headheight}{15pt}

% Estilo de páginas iniciales
\fancypagestyle{plain}{
    \fancyhf{}
    \fancyfoot[C]{\small\thepage}
    \renewcommand{\headrulewidth}{0pt}
}

% Formato de Títulos
\usepackage{titlesec}
\titleformat{\section}
    {\normalfont\large\bfseries}{\thesection.}{1em}{}
\titleformat{\subsection}
    {\normalfont\normalsize\bfseries}{\thesubsection.}{1em}{}
\titleformat{\subsubsection}
    {\normalfont\normalsize\itshape}{\thesubsubsection.}{1em}{}

% Control de listas numeradas y con viñetas
\usepackage{enumitem}
\setlist[itemize]{noitemsep, topsep=2pt, leftmargin=1.5cm}
\setlist[enumerate]{noitemsep, topsep=2pt, leftmargin=1.5cm}

% Compatibilidad con Pandoc
\providecommand{\tightlist}{%
  \setlength{\itemsep}{0pt}\setlength{\parskip}{0pt}}

% Colores y Enlaces (Modo Tinta Negra Oficial para Tesis y Protocolos UNSCH)
\usepackage{xcolor}

\usepackage{hyperref}
\hypersetup{
    colorlinks=true,
    linkcolor=black,
    citecolor=black,
    urlcolor=black,
    filecolor=black,
    pdfauthor={Gresly Lucero Pariona Palomino},
    pdftitle={Acceso a los servicios de Salud y Calidad de Atencion percibido por los usuarios},
    pdfsubject={Proyecto de Investigacion en Salud (EN 486) - UNSCH},
    pdfkeywords={Acceso a los servicios de salud, Calidad de atencion, Percepcion del usuario, Satisfaccion, Enfermeria, UNSCH, Ayacucho}
}

% Comillas en español
\usepackage{csquotes}

% Microtipografía
\usepackage{microtype}

% Compatibilidad con bloques de codigo de Pandoc
\ifx\Shaded\undefined
  \newenvironment{Shaded}{\begin{quote}\small\ttfamily}{\end{quote}}
\fi
\ifx\Highlighting\undefined
  \newenvironment{Highlighting}{}{}
\fi


\begin{document}
\section{INVESTIGACIÓN CIENTÍFICA}\label{investigaciuxf3n-cientuxedfica}  \textbf{UNIVERSIDAD NACIONAL DE SAN CRISTÓBAL DE HUAMANGA}\\ \textbf{FACULTAD DE CIENCIAS DE LA SALUD}\\ \textbf{ESCUELA PROFESIONAL DE ENFERMERÍA}  \begin{itemize} \tightlist \item   \textbf{Asignatura:} Proyecto de Investigación en Salud (EN 486)\\ \item   \textbf{Docente:} Dr.~Manglio Aguirre Andrade    \begin{itemize}   \tightlist   \item     \emph{Docente en la Categoría Principal de la UNSCH}\\   \item     \emph{Experto en Salud Pública}\\   \item     \textbf{Correo:} manglio.aguirre@unsch.edu.pe\\   \item     \textbf{Teléfono:} 966102220   \end{itemize} \end{itemize}  \begin{center}\rule{0.5\linewidth}{0.5pt}\end{center}  \subsection{1. INVESTIGACIÓN}\label{investigaciuxf3n}  \subsubsection{Etimología}\label{etimologuxeda}  \begin{itemize} \tightlist \item   \textbf{In:} en, hacia \item   \textbf{Vestigium:} huella, pista \item   \emph{Investigar es caminar hacia o hacer seguimiento de huellas o   pistas.} \end{itemize}  \subsubsection{Concepto}\label{concepto}  \begin{itemize} \tightlist \item   El objetivo de la investigación es descubrir algo nuevo y para ello se   requiere la aplicación formal de métodos. \item   ¿Investigar es equivalente a\ldots? \end{itemize}  \begin{center}\rule{0.5\linewidth}{0.5pt}\end{center}  \subsection{1. INVESTIGACIÓN CIENTÍFICA}\label{investigaciuxf3n-cientuxedfica-1}  \begin{quote} \textbf{Neil Salkind:}\\ La investigación es un proceso sistemático mediante el cual se descubren conocimientos nuevos. \end{quote}  \subsubsection{Características de la Investigación Científica:}\label{caracteruxedsticas-de-la-investigaciuxf3n-cientuxedfica}  \begin{itemize} \tightlist \item   \textbf{Se basa en el trabajo de otros:} Se articula sobre el   conocimiento científico acumulado. \item   \textbf{Se puede repetir (replicabilidad):} Otros investigadores   pueden reproducir los procedimientos y verificar los resultados. \item   \textbf{Se puede generalizar a otras situaciones:} Sus conclusiones   trascienden los casos particulares estudiados. \item   \textbf{Se basa en algún razonamiento lógico y está vinculado a una   teoría:} Integra marcos conceptuales coherentes. \item   \textbf{Se puede hacer (viabilidad/factibilidad):} Es ejecutable con   los recursos, tiempos y técnicas disponibles. \item   \textbf{Genera nuevas preguntas, es incremental:} No es estática; abre   nuevos campos de indagación. \item   \textbf{Se sustenta en el método científico:} Emplea procedimientos   rigurosos de contrastación. \end{itemize}  \begin{center}\rule{0.5\linewidth}{0.5pt}\end{center}  \subsection{2. MÉTODO CIENTÍFICO}\label{muxe9todo-cientuxedfico}  \subsubsection{Definición}\label{definiciuxf3n}  \begin{quote} \textbf{Francis Bacon:}\\ El método científico es un conjunto de reglas para observar fenómenos e inferir conclusiones a partir de dichas observaciones. \end{quote}  Es un procedimiento riguroso formado por una \textbf{secuencia lógica de actividades} que procura descubrir: 1. Las \textbf{características} de los fenómenos. 2. Las \textbf{relaciones internas} entre sus elementos constitutivos. 3. Sus \textbf{conexiones} con otros fenómenos de la realidad.  \emph{La secuencia lógica está guiada por el raciocinio y la comprobación a través de la demostración y la verificación.}  \begin{center}\rule{0.5\linewidth}{0.5pt}\end{center}  \subsubsection{Veracidad y Verificabilidad en la Ciencia}\label{veracidad-y-verificabilidad-en-la-ciencia}  Una condición fundamental de la ciencia es la veracidad y la verificabilidad. Los enunciados verificables son de diversas clases:  \begin{enumerate} \def\labelenumi{\arabic{enumi}.} \tightlist \item   \textbf{Proposiciones singulares:}\\   \emph{Ejemplo:} \emph{``Este trozo de hierro está caliente''}. \item   \textbf{Proposiciones particulares o existenciales:}\\   \emph{Ejemplo:} \emph{``Algunos trozos de hierro están calientes''}   (proposición verificablemente falsa en contextos determinados). \item   \textbf{Enunciados de leyes:}\\   \emph{Ejemplo:} \emph{``Todos los metales se dilatan con el calor''}   (formalmente: \emph{``Para todo \(x\), si \(x\) es un trozo de metal   que se calienta, entonces \(x\) se dilata''}). \end{enumerate}  \begin{quote} \textbf{Niveles de verificación:}\\ * Las proposiciones singulares y particulares pueden verificarse a menudo de manera inmediata, con el auxilio directo de los sentidos o de instrumentos que amplían su alcance.\\ * Otras veces exigen operaciones complejas que implican leyes científicas y cálculos matemáticos (por ejemplo: la medición de distancias astronómicas como la distancia media Tierra-Sol). \end{quote}  \begin{center}\rule{0.5\linewidth}{0.5pt}\end{center}  \subsubsection{Modo de Conocer y Proceder de la Ciencia}\label{modo-de-conocer-y-proceder-de-la-ciencia}  \begin{verbatim}     ┌──────────────────────────┐     │  PROBLEMA / IDEA /       │     │  OBSTÁCULO               │     └────────────┬─────────────┘                  │                  ▼     ┌──────────────────────────┐     │  HIPÓTESIS               │     └────────────┬─────────────┘                  │                  ▼     ┌──────────────────────────┐     │  RAZONAMIENTO -          │     │  DEDUCCIÓN               │     └────────────┬─────────────┘                  │                  ▼     ┌──────────────────────────┐     │  OBSERVACIÓN, PRUEBA,    │     │  EXPERIMENTO             │     └────────────┬─────────────┘                  │                  ▼     ┌──────────────────────────┐     │  PRODUCCIÓN DE CIENCIA   │     └──────────────────────────┘ \end{verbatim}  El método científico se basa en: 1. La observación sistemática de la realidad. 2. Su medición cuantitativa y cualitativa. 3. El análisis riguroso de sus propiedades y características. 4. La elaboración formal de hipótesis explicativas. 5. Su interpretación y contrastación empírica. 6. La formulación de alternativas de acción o respuesta para transformar la realidad.  \begin{center}\rule{0.5\linewidth}{0.5pt}\end{center}  \subsection{3. TIPOS DE INVESTIGACIÓN}\label{tipos-de-investigaciuxf3n}  \emph{(Ref. Landeau Rebeca, 2007)}  \subsubsection{a. Según la Finalidad}\label{a.-seguxfan-la-finalidad}  \begin{itemize} \tightlist \item   \textbf{a.1. Investigación teórica, básica o pura:}\\   Se interesa en el descubrimiento de las leyes que rigen el   comportamiento de ciertos fenómenos o eventos; intenta encontrar los   principios generales que gobiernan los diversos fenómenos en los que   el investigador se encuentra interesado \emph{(Kerlinger, 1982)}. \item   \textbf{a.2. Investigación aplicada:}\\   Aquella que incluye cualquier esfuerzo sistemático y socializado por   resolver problemas o intervenir situaciones, aunque no sea   programático, es decir, aunque no pertenezca a una trayectoria previa   de investigaciones descriptivas y teóricas. Es aplicable tanto a la   innovación técnica, artesanal e industrial como a la propiamente   científica \emph{(Padrón, 2006)}. \end{itemize}  \subsubsection{b. Según su Carácter}\label{b.-seguxfan-su-caruxe1cter}  \begin{itemize} \tightlist \item   \textbf{b.1. Investigación exploratoria:} Primer acercamiento a temas   poco estudiados o desconocidos. \item   \textbf{b.2. Investigación descriptiva:} Caracteriza situaciones,   eventos, propiedades o perfiles de personas o comunidades. \item   \textbf{b.3. Investigación correlacional o \emph{ex post facto}:} Mide   el grado de relación o asociación entre dos o más variables sin   manipulación directa. \item   \textbf{b.4. Investigación explicativa:} Dirigida a responder por las   causas de los eventos y fenómenos físicos o sociales. \item   \textbf{b.5. Investigación experimental:} Manipulación deliberada de   variables independientes para observar efectos en variables   dependientes bajo estricto control. \end{itemize}  \subsubsection{c.~Según su Naturaleza}\label{c.-seguxfan-su-naturaleza}  \begin{itemize} \tightlist \item   \textbf{c.1. Investigación cuantitativa:} Recolección y análisis de   datos cuantificables, medición numérica y contrastación estadística. \item   \textbf{c.2. Investigación cualitativa:} Comprensión profunda de   significados, vivencias, discursos y fenómenos en su contexto natural. \item   \textbf{c.3. Investigación cuali-cuantitativa (Mixta):} Integración   sistemática de enfoques cuantitativos y cualitativos. \end{itemize}  \subsubsection{d.~Según su Alcance Temporal}\label{d.-seguxfan-su-alcance-temporal}  \begin{itemize} \tightlist \item   \textbf{d.1. Investigación transversal (sincrónica):} Recolección de   datos en un solo momento o punto temporal único. \item   \textbf{d.2. Investigación longitudinal (diacrónica):} Recolección de   datos a través del tiempo en períodos sucesivos para analizar   evoluciones o tendencias. \end{itemize}  \subsubsection{e. Según la Orientación que Asume}\label{e.-seguxfan-la-orientaciuxf3n-que-asume}  \begin{itemize} \tightlist \item   \textbf{e.1. Investigación orientada a la comprobación:} Dirigida a   contrastar o validar hipótesis, teorías o leyes. \item   \textbf{e.2. Investigación orientada al descubrimiento:} Enfocada en   generar nuevas ideas, patrones y hallazgos teóricos. \item   \textbf{e.3. Investigación orientada a la aplicación:} Focalizada en   dar solución inmediata o práctica a problemas reales. \end{itemize}  \begin{center}\rule{0.5\linewidth}{0.5pt}\end{center}  \subsection{4. ETAPAS DEL PROCESO DE INVESTIGACIÓN}\label{etapas-del-proceso-de-investigaciuxf3n}  \begin{Shaded} \begin{Highlighting}[] \NormalTok{flowchart TD} \NormalTok{    A["1. Elaboración del Proyecto de Investigación"] {-}{-}\textgreater{} B["APROBACIÓN"]} \NormalTok{    B {-}{-}\textgreater{} C["2. Ejecución del Proyecto de Investigación"]} \NormalTok{    C {-}{-}\textgreater{} D["3. Elaboración del Informe Final o Tesis"]} \NormalTok{    D {-}{-}\textgreater{} E["APROBACIÓN"]} \NormalTok{    E {-}{-}\textgreater{} F["4. Sustentación de la Tesis"]} \NormalTok{    F {-}{-}\textgreater{} G["5. Publicación"]} \end{Highlighting} \end{Shaded}  \begin{center}\rule{0.5\linewidth}{0.5pt}\end{center}  \subsection{5. COMPROMISOS SOBRE LA BIOÉTICA Y ÉTICA DE LA INVESTIGACIÓN}\label{compromisos-sobre-la-biouxe9tica-y-uxe9tica-de-la-investigaciuxf3n}  \begin{itemize} \tightlist \item   \textbf{¿Cómo encontrar la verdad?\ldots{} ¿Es la verdad?} \item   \textbf{¿Explorando? ¿Examinando? ¿Vulnerando? ¿Manipulando?} \item   \textbf{Rol de los Comités de Bioética:} Garantizar la protección de   los derechos, bienestar e integridad de los sujetos participantes, el   respeto irrestricto de la confidencialidad, el consentimiento   informado y el cumplimiento de principios de beneficencia, no   maleficencia, autonomía y justicia. \end{itemize}  \begin{center}\rule{0.5\linewidth}{0.5pt}\end{center}  \subsection{6. ESQUEMAS DE PROTOCOLOS DE INVESTIGACIÓN}\label{esquemas-de-protocolos-de-investigaciuxf3n}  \subsubsection{A. Esquema o Protocolo de un Proyecto de Investigación con Enfoque Cuantitativo}\label{a.-esquema-o-protocolo-de-un-proyecto-de-investigaciuxf3n-con-enfoque-cuantitativo}  \paragraph{I. EL PROBLEMA DE INVESTIGACIÓN}\label{i.-el-problema-de-investigaciuxf3n}  \begin{itemize} \tightlist \item   \textbf{1.1.} Origen del problema de investigación \item   \textbf{1.2.} Delimitación del problema \item   \textbf{1.3.} Formulación del problema \item   \textbf{1.4.} Objetivos    \begin{itemize}   \tightlist   \item     Objetivo general   \item     Objetivos específicos   \end{itemize} \item   \textbf{1.5.} Justificación del estudio \end{itemize}  \paragraph{II. MARCO TEÓRICO}\label{ii.-marco-teuxf3rico}  \begin{itemize} \tightlist \item   \textbf{2.1.} Antecedentes del estudio \item   \textbf{2.2.} Base teórico-científica \item   \textbf{2.3.} Hipótesis \item   \textbf{2.4.} Identificación y operacionalización de variables \end{itemize}  \paragraph{III. METODOLOGÍA DE INVESTIGACIÓN}\label{iii.-metodologuxeda-de-investigaciuxf3n}  \begin{itemize} \tightlist \item   \textbf{3.1.} Tipo de investigación \item   \textbf{3.2.} Tipo de diseño de investigación \item   \textbf{3.3.} Área de estudio \item   \textbf{3.4.} Población \item   \textbf{3.5.} Muestra \item   \textbf{3.6.} Técnica e instrumento de recolección de datos \item   \textbf{3.7.} Plan de recolección de datos \item   \textbf{3.8.} Plan de procesamiento de datos \item   \textbf{3.9.} Plan de presentación y análisis de datos \end{itemize}  \paragraph{IV. ASPECTOS ADMINISTRATIVOS}\label{iv.-aspectos-administrativos}  \begin{itemize} \tightlist \item   \textbf{4.1.} Asignación de recursos \item   \textbf{4.2.} Cronograma \item   \textbf{4.3.} Financiamiento \end{itemize}  \paragraph{REFERENCIAS BIBLIOGRÁFICAS}\label{referencias-bibliogruxe1ficas}  \paragraph{ANEXOS:}\label{anexos}  \begin{itemize} \tightlist \item   Instrumento de recolección de datos. \item   Matriz de consistencia. \item   Validación del instrumento por juicio de expertos, consentimiento   informado / autoinformado. \end{itemize}  \begin{center}\rule{0.5\linewidth}{0.5pt}\end{center}  \subsubsection{B. Esquema o Protocolo de un Proyecto de Investigación con Enfoque Cualitativo}\label{b.-esquema-o-protocolo-de-un-proyecto-de-investigaciuxf3n-con-enfoque-cualitativo}  \paragraph{INTRODUCCIÓN}\label{introducciuxf3n}  \paragraph{CAPÍTULO I: EL PROBLEMA}\label{capuxedtulo-i-el-problema}  \begin{itemize} \tightlist \item   Problema de investigación \item   Preguntas norteadoras \item   Planteamiento del problema central o formulación del enunciado del   problema de investigación \item   Objetivos \item   Justificación \end{itemize}  \paragraph{CAPÍTULO II: MARCO CONTEXTUAL}\label{capuxedtulo-ii-marco-contextual}  \paragraph{CAPÍTULO III: ABORDAJE TEÓRICO}\label{capuxedtulo-iii-abordaje-teuxf3rico}  \begin{itemize} \tightlist \item   Antecedentes de estudio \item   Descripción de las variables / categorías de estudio\\   \emph{(Puede plantearse una hipótesis de trabajo)} \end{itemize}  \paragraph{CAPÍTULO IV: METODOLOGÍA CUALITATIVA DEL ESTUDIO}\label{capuxedtulo-iv-metodologuxeda-cualitativa-del-estudio}  \emph{(Deberá seleccionarse tanto el método general como el específico)} * \textbf{Método General:} * Estudios descriptivos (Demográficos, estudios de casos) * Estudios interpretativos * Estudios históricos * \textbf{Métodos Cualitativos Específicos:} * Método hermenéutico * Método dialéctico * Método fenomenológico * Método etnográfico * Método de investigación-acción (I-A) * Método de historias de vida * Método dialéctico estructural * \textbf{Procedimientos metodológicos:} * Selección de los informantes / muestra cualitativa * Técnicas e instrumentos de recolección de información * Plan de recolección de datos * Plan de procesamiento y categorización de datos * Plan de presentación y análisis e interpretación de resultados  \paragraph{CAPÍTULO V: ASPECTOS ADMINISTRATIVOS}\label{capuxedtulo-v-aspectos-administrativos}  \begin{itemize} \tightlist \item   Asignación de recursos \item   Cronograma \item   Financiamiento \end{itemize}  \paragraph{BIBLIOGRAFÍA}\label{bibliografuxeda}  \paragraph{ANEXOS:}\label{anexos-1}  \begin{itemize} \tightlist \item   Guía de entrevista / instrumentos de recolección cualitativa. \item   Matriz de consistencia cualitativa. \end{itemize}  \begin{center}\rule{0.5\linewidth}{0.5pt}\end{center}  \emph{FACULTAD DE CIENCIAS DE LA SALUD - ESCUELA PROFESIONAL DE ENFERMERÍA}\\ \emph{Gracias.}
\end{document}
